% References: https://eel.is/c++draft/alg.sorting
% https://en.cppreference.com/w/cpp/algorithm/nth_element.html
\subsection{排序、选第 k 小与区间操作}
\label{stl:sort}
\subsubsection*{原理}
\InlineCode{sort} 对随机访问区间排序，比较次数为 $O(n\log n)$；比较器应表示“左项应该排在右项前面”。规则必须严格且一致：相等时返回假，不能用小于等于，也不要用可能溢出的两数相减来比较。
\subsubsection*{使用教程}
\InlineCode{sort(first,last)} 默认升序，加 \InlineCode{greater<int>()} 则降序；\InlineCode{stable_sort} 保留等价项原有相对顺序。\InlineCode{nth_element(first,nth,last)} 平均线性地把该排名应在的值放到 \InlineCode{nth}，左右两侧不保证各自有序。
\begin{lstlisting}[style=mtycpp,numbers=none]
vector<int> a={7,2,5,2,9};
int k=3;
auto b=a;
nth_element(b.begin(),b.begin()+k-1,b.end());
int kth=b[k-1];
sort(a.begin(),a.end(),greater<int>());
vector<pair<int,int>> p={{2,3},{1,3},{4,1}};
sort(p.begin(),p.end(),[](const auto &x,const auto &y){
	return tie(x.second,x.first)<tie(y.second,y.first);
});
\end{lstlisting}
第 $k$ 小按 $1\le k\le n$ 使用偏移 $k-1$。\InlineCode{reverse(first,last)} 反转区间，\InlineCode{find(first,last,x)} 查首次出现，\InlineCode{min_element/max_element} 返回极值位置，均为线性；查找或极值在空区间返回尾迭代器，不能直接解引用。
\begin{lstlisting}[style=mtycpp,numbers=none]
vector<int> a={3,1,4};
reverse(a.begin(),a.end());
auto it=max_element(a.begin(),a.end());
if(it!=a.end()){int best=*it;}
long long x=min({3LL,5LL,2LL});
long long y=clamp(x,1LL,4LL);
\end{lstlisting}
\InlineCode{clamp(x,l,r)} 将值限制到闭区间，要求 $l\le r$；同一次 \InlineCode{min/max/clamp} 调用的参数类型须一致，整数宏不会自动改变字面量的类型，64 位常量按需加 \InlineCode{LL}。占位下标不参与算法：对 1 下标数组使用 \InlineCode{begin()+1}。
